% -*- coding: utf-8 -*-
% !TEX program = xelatex
\documentclass[plain,noanswer,medmath]{../../template/jnuexam}
\usepackage{../../template/kysx}

\begin{document}


\examtitle{2000}{1}


\makepart{填空题}{本题共5小题，每小题3分，满分15分}


\begin{problem}
$\int_0^1 \sqrt{2x - x^2} \dx =$ \fillin{}．
\end{problem}


\begin{problem}
曲面$x^2 + 2y^2 + 3z^2 = 21$在点$\left(1,\text{-}2,2 \right)$的法线方程为 \fillin{}．
\end{problem}


\begin{problem}
微分方程$xy'' + 3y' = 0$的通解为 \fillin{}．
\end{problem}


\begin{problem}
已知方程组$\begin{pmatrix}
  1 & 2 & 1 \\
  2 & 3 & a + 2 \\
  1 & a & -2
\end{pmatrix}\begin{pmatrix}
  x_1 \\
  x_2 \\
  x_3
\end{pmatrix} = \begin{pmatrix}
  1 \\
  3 \\
  0
\end{pmatrix}$无解，则$a =$ \fillin{}．
\end{problem}


\begin{problem}
设两个相互独立的事件$A$和$B$都不发生的概率为$\frac{1}{9}$ , $A$发生$B$不发生的概率与$B$发生$A$不发生的概率相等，
则$P(A)$ = \fillin{}．
\end{problem}


\makepart{选择题}{本题共5小题，每小题3分，共15分}


\begin{problem}
设$f(x),g(x)$是恒大于零的可导函数，且$f'(x)g(x) - f(x)g'(x) < 0$，则当$a < x < b$ 时，有\pickout{}
\begin{abcd}
\item $f(x)g(b) > f(b)g(x)$．
\item $f(x)g(a) > f(a)g(x)$．
\item $f(x)g(x) > f(b)g(b)$．
\item $f(x)g(x) > f(a)g(a)$．
\end{abcd}
\end{problem}


\begin{problem}
设$S: x^2 + y^2 + z^2 = a^2\;(z \ge 0)$，$S_1$为$S$在第一卦限中的部分，则有\pickout{}
\begin{abcd}
\item $\iint_S x\dS = 4\iint_{S_1} x\dS$．
\item $\iint_S y\dS = 4\iint_{S_1} x\dS$．
\item $\iint_S z\dS = 4\iint_{S_1} x\dS$．
\item $\iint_S xyz\dS = 4\iint_{S_1} xyz\dS$．
\end{abcd}
\end{problem}


\begin{problem}
设级数$\sum_{n = 1}^{\infty}u_n$收敛，则必收敛的级数为\pickout{}
\begin{abcd}
\item $\sum_{n = 1}^{\infty}(- 1)^n\frac{u_n}{n}$．
\item $\sum_{n = 1}^{\infty}u_n^2$．
\item $\sum_{n = 1}^{\infty}(u_{2n - 1} - u_{2n})$．
\item $\sum_{n = 1}^{\infty}(u_n + u_{n + 1})$．
\end{abcd}
\end{problem}


\begin{problem}
设$n$维列向量组$\alpha_1, \cdot \cdot \cdot ,\alpha_m(m < n)$线性无关，则$n$维列向量组$\beta_1, \cdot \cdot \cdot ,\beta_m$线性无关的充分必要条件为\pickout{}
\begin{abcd}
\item 向量组$\alpha_1, \cdot \cdot \cdot ,\alpha_m$可由向量组$\beta_1, \cdot \cdot \cdot ,\beta_m$线性表示．
\item 向量组$\beta_1, \cdot \cdot \cdot ,\beta_m$可由向量组$\alpha_1, \cdot \cdot \cdot ,\alpha_m$线性表示．
\item 向量组$\alpha_1, \cdot \cdot \cdot ,\alpha_m$与向量组$\beta_1, \cdot \cdot \cdot ,\beta_m$等价．
\item 矩阵$A = \left(\alpha_1, \cdot \cdot \cdot ,\alpha_m \right)$与矩阵$B = \left(\beta_1, \cdot \cdot \cdot ,\beta_m \right)$等价．
\end{abcd}
\end{problem}


\begin{problem}
设二维随机变量$\left(X,Y \right)$服从二维正态分布，则随机变量$\xi = X + Y$与$\eta = X - Y$不相关的充分必要条件为\pickout{}
\begin{abcd}
\item $E(X) = E(Y).$
\item $E(X^2) - \left[E(X) \right]^2 = E(Y^2) - \left[E(Y) \right]^2.$
\item $E(X^2) = E(Y^2).$
\item $E(X^2) + \left[E(X) \right]^2 = E(Y^2) + \left[E(Y) \right]^2.$
\end{abcd}
\end{problem}


\makepart{}{本题满分5分}

\begin{problem}
求$\lim_{x \to 0} \left(\frac{2 + \e^{\frac{1}{x}}}{1 + \e^{\frac{4}{x}}} + \frac{\sin x}{|x|} \right).$
\end{problem}


\makepart{}{本题满分6分}

\begin{problem}
设$z = f\left(xy,\frac{x}{y} \right) + g\left(\frac{x}{y} \right)$，
其中$f$具有二阶连续偏导数，$g$具有二阶连续导数，求$\frac{\pd^2 z}{\pdx\pdy}.$
\end{problem}


\makepart{}{本题满分6分}

\begin{problem}
计算曲线积分$I = \oint_L \frac{x\dy - y\dx}{4x^2 + y^2} $，
其中$L$是以点$\left(1,0 \right)$为中心，$R$为半径的圆周$(R > 1)$，取逆时针方向．
\end{problem}


\makepart{}{本题满分7分}

\begin{problem}
设对于半空间$x > 0$内任意的光滑有向封闭曲面$S$，都有
$$\oiint_S xf(x) \dy\dz - xyf(x)\dz\dx - \e^{2x} z\dx\dy = 0,$$
其中函数$f(x)$在$(\text{0, +}\infty)$内具有连续的一阶导数，且$\lim_{x \to 0^+} f(x) = 1$，求$f(x)$．
\end{problem}


\makepart{}{本题满分6分}

\begin{problem}
求幂级数$\sum_{n = 1}^{\infty}\frac{1}{3^n + (- 2)^n} \frac{x^n}{n}$的收敛区域，并讨论该区间端点处的收敛性．
\end{problem}


\makepart{}{本题满分7分}

\begin{problem}
设有一半径为$R$的球体，$P_0$是此球的表面上的一个定点，
球体上任一点的密度与该点到$P_0$距离的平方成正比(比例常数$k > 0$)，求球体的重心位置．
\end{problem}


\makepart{}{本题满分6分}

\begin{problem}
设函数$f(x)$在$\left[0,\pi \right]$上连续，且$\int_0^{\pi}f(x)\dx = 0$，$\int_0^{\pi}f(x)\cos x\dx = 0$．
试证：在$(0,\pi)$内至少存在两个不同的点$\xi_1,\xi_2$，使$f(\xi_1) = f(\xi_2) = 0$．
\end{problem}


\makepart{}{本题满分6 分}

\begin{problem}
设矩阵$A$的伴随矩阵$A^* = \begin{pmatrix}
  1 &  0 & 0 & 0 \\
  0 &  1 & 0 & 0 \\
  1 &  0 & 1 & 0 \\
  0 & -3 & 0 & 8
\end{pmatrix}$，且$AB A^{-1} = BA^{-1} + 3E$，其中$E$为$4$阶单位矩阵，求矩阵$B$．
\end{problem}


\makepart{}{本题满分8分}

\begin{problem}
某试验性生产线每年一月份进行熟练工与非熟练工的人数统计，
然后将$\frac{1}{6}$熟练工支援其他生产部门，其缺额由招收新的非熟练工补齐，
新、老非熟练工经过培训及实践至年终考核有$\frac{2}{5}$成为熟练工．
设第$n$年一月份统计的熟练工和非熟练工所占百分比分别为$x_n,$$y_n$记成向量$\begin{pmatrix}
  x_n \\
  y_n
\end{pmatrix}$．
\begin{enumerate}
\item 求$\begin{pmatrix}
  x_{n + 1} \\
  y_{n + 1}
\end{pmatrix}$与$\begin{pmatrix}
  x_n \\
  y_n
\end{pmatrix}$的关系式并写成矩阵形式：$\begin{pmatrix}
  x_{n + 1} \\
  y_{n + 1}
\end{pmatrix} = A\begin{pmatrix}
  x_n \\
  y_n
\end{pmatrix};$
\item 验证$\eta_1 = \begin{pmatrix}
  4 \\
  1
\end{pmatrix},\eta_2 = \begin{pmatrix}
  -1 \\
  1
\end{pmatrix}$是$A$的两个线性无关的特征向量，并求出相应的特征值；
\item 当$\begin{pmatrix}
  x_1 \\
  y_1
\end{pmatrix} = \begin{pmatrix}
  \frac{1}{2} \\
  \frac{1}{2}
\end{pmatrix}$时，求$\begin{pmatrix}
  x_{n + 1} \\
  y_{n + 1}
\end{pmatrix}.$
\end{enumerate}
\end{problem}


\makepart{}{本题满分8分}

\begin{problem}
某流水生产线上每个产品不合格的概率为$p$（$0 < p < 1$），各产品合格与否相互独立，
当出现一个不合格产品时即停机检修．设开机后第一次停机时已生产了产品的个数为$X$，求$X$的数学期望$E(X)$和方差$D(X)$．
\end{problem}


\makepart{}{本题满分8分}

\begin{problem}
设某种元件的使用寿命$X$的概率密度为$f(x;\theta) = \begin{cases}
  2\e^{- 2(x - \theta)},x \ge \theta \\
  0, x < \theta
\end{cases}$ 其中$\theta > 0$为未知参数，又设$x_1,x_2, \cdot \cdot \cdot ,x_n$是$X$的一组样本观测值，
求参数$\theta$的最大似然估计值．
\end{problem}


\end{document}
